% Genere par LIVRABLES/PHASE2/profilage_duckdb.R
\section{Registre des problemes de qualite}

{\footnotesize\setlength{\tabcolsep}{3pt}
\begin{longtable}{@{}p{1.1cm}p{2.2cm}p{4.2cm}p{3.6cm}p{4.3cm}@{}}
\toprule
\textbf{Id} & \textbf{Categorie} & \textbf{Probleme} & \textbf{Quantification} & \textbf{Action corrective} \\
\midrule \endfirsthead
\toprule \textbf{Id} & \textbf{Categorie} & \textbf{Probleme} & \textbf{Quantification} & \textbf{Action corrective} \\ \midrule \endhead
\bottomrule \endfoot
\textbf{PQ-01} & Grain / duplication & DIS\_PLV est eclate par reseau amont : ce n'est pas le grain du prelevement & 4 060 859 lignes pour 2 834 678 prelevements distincts, soit +43,3 \% de lignes. Un prelevement apparait jusqu'a 66 fois & Dedoublonner sur referenceprel avant de charger le fait PRELEVEMENT. Controle prealable valide : 0 prelevement rattache a plusieurs communes, 0 prelevement avec des conformites divergentes \\[2pt]
\textbf{PQ-02} & Integrite referentielle & Le reseau amont reference des codes absents du referentiel des UDI & 1 114 396 lignes (27,4 \%) et 21 139 codes distincts sans correspondance dans DIS\_COM\_UDI & Table pont RESEAU\_AMONT avec un membre Inconnu. Ne pas utiliser le reseau amont comme niveau hierarchique fiable \\[2pt]
\textbf{PQ-03} & Format & La part de debit est stockee en texte avec le signe pourcent & 56,04 \% de valeurs absentes, 101 modalites de 0 \% a 100 \%, dont 25 477 lignes a 0 \% & Retirer le signe et caster en decimal. Definir la regle applicable aux poids nuls \\[2pt]
\textbf{PQ-04} & Format & Heure de prelevement au format non respecte & 5 880 valeurs nulles et des formats corrompus observes : 9:h54, 8:h30, 10h, 00h0, h & Parser avec expression reguliere et table de rejet. Attribut optionnel : la dimension temps au jour suffit au workload \\[2pt]
\textbf{PQ-05} & Typage & Les codes geographiques perdent leur signification s'ils sont lus en numerique & Formats 001 a 976 plus 02A et 02B cote SISE-Eaux, 01001 cote INSEE & Forcer la lecture en VARCHAR (all\_varchar=true dans DuckDB, type String dans KNIME) \\[2pt]
\textbf{PQ-06} & Valeurs sentinelles & Des valeurs conventionnelles remplacent l'absence de donnee & siren\_epci = NR sur 10 lignes, insee\_can = NR sur 119, statut nul sur 12 865, date sentinelle 1900-01-01 & Mapper NR et 1900-01-01 vers NULL a l'import \\[2pt]
\textbf{PQ-07} & Semantique & La commune du prelevement n'est pas la commune desservie & 34 915 communes principales pour 34 944 communes desservies. Relation reelle N-N via DIS\_COM\_UDI (495 059 liens) & Decision de modelisation a expliciter : arc multiple avec table pont, ou rattachement simplifie a la commune principale. Documenter le biais retenu \\[2pt]
\textbf{PQ-08} & Texte libre & La conclusion sanitaire est un texte redige sans nomenclature & 140 609 variantes distinctes, longueur jusqu'a 3 782 caracteres & Classifier en 5 a 8 modalites par motifs textuels pour alimenter DIM\_CONCLUSION. Conserver le texte source en attribut non analytique \\[2pt]
\textbf{PQ-09} & Qualite des libelles & Les libelles geographiques ne sont pas des identifiants & 32 694 noms distincts pour 34 944 communes : homonymes. Colonne quartier non normalisee (17 183 modalites, valeurs du type (1\%) ou epizy, centre ville) & Joindre exclusivement sur les codes INSEE. Rejeter la colonne quartier \\[2pt]
\textbf{PQ-10} & Normalisation & Libelles non trimes, casse incoherente, libelle variable pour un meme code & 44 560 valeurs de ugelib avec espaces ou tabulations parasites, 25 670 conclusions non trimees, 1 914 libelles pour 1 767 codes parametre, 72 unites avec doublons de casse, 25 849 noms de reseau pour 24 373 codes & Trim, normalisation de casse, et pour chaque code retenir le libelle le plus recent (SCD type 1) \\[2pt]
\textbf{PQ-11} & Codification & Les indicateurs de conformite ont plus de deux modalites & Bacterio : C 3 715 800, S 287 850, N 57 153. Chimique : C 3 952 830, N 101 985, D 3 906, S 2 082 & Mesure = 1 si N, 0 si C, NULL sinon. Exclure S du denominateur. Statuer explicitement sur D (derogation) et tracer la decision dans le rapport \\[2pt]
\textbf{PQ-12} & Colonnes inexploitables & Colonnes vides, redondantes ou purement techniques & libwebparametre : 921 valeurs non vides sur 125 167 285 lignes. referenceanl absent sur 11,73 \%. cdparametre nul sur 0,10 \% & Rejet documente. Cle du fait ANALYSE = referenceprel + cdparametresiseeaux et non referenceanl \\[2pt]
\textbf{PQ-13} & Heterogeneite & Des parametres qualitatifs sont melanges aux parametres quantitatifs & 9 791 250 lignes (7,82 \%) avec qualitparam = O (odeur, aspect, saveur) & Filtrer sur qualitparam = N pour toute mesure numerique. Conserver les qualitatifs pour les comptages \\[2pt]
\textbf{PQ-14} & Heterogeneite & Le resultat brut melange nombres, valeurs censurees et texte & 48 773 modalites distinctes. Aucun changement anormal 6 881 300, Aspect normal 2 422 099, N.M. 919 918, SEUIL 220 843, valeurs negatives sentinelles -1,0 / -2,0 / -3,0, et le caractere ? sur 1 228 lignes & Ne jamais caster rqana. Utiliser valtraduite pour la mesure et rqana uniquement pour deriver les indicateurs de censure \\[2pt]
\textbf{PQ-15} & Format & Les seuils reglementaires sont du texte libre avec operateur et unite & limitequal absent sur 43,70 \%, refqual sur 80,28 \%. 64 118 644 seuils utilisent la virgule decimale contre 1 681 966 le point. 6 445 563 seuils sont des intervalles a deux bornes du type superieur ou egal a 6,5 et inferieur ou egal a 9 & Parser en quatre champs (operateur, borne\_min, borne\_max, unite) pour alimenter DIM\_SEUIL. Normaliser le separateur decimal avant cast \\[2pt]
\textbf{PQ-16} & Censure statistique & Les resultats inferieurs a la limite de quantification sont codes a zero & 82 653 622 resultats censures a gauche (66,03 \% du total), dont 82 651 485 avec valtraduite exactement egale a 0, soit 99,997 \%. 76 478 censures a droite & Ajouter un indicateur est\_censure. Pour les analyses de concentration, appliquer une substitution documentee (LQ/2) ou restreindre aux quantifications. Privilegier le taux de depassement plutot que la moyenne \\[2pt]
\textbf{PQ-17} & Couverture temporelle & Le millesime 2021 de Solagro est partiel & 19 587 communes en 2021 contre 34 806 en 2020 et 2022 & Exclure 2021 des comparaisons ou le declarer explicitement comme partiel dans les visualisations \\[2pt]
\textbf{PQ-18} & Integration & Le code departement n'a pas le meme format dans les deux sources & 3 caracteres cote SISE-Eaux (001) contre 2 cote Solagro (01) & Normaliser sur un format unique au niveau de DIM\_GEOGRAPHIE. Traiter 2A et 2B separement \\[2pt]
\textbf{PQ-19} & Hierarchie variable & Deux decoupages regionaux coexistent dans la source & 12 codes dans le decoupage post-2016 contre 20 dans l'ancien & Retenir le decoupage post-2016 comme niveau du DW. Conserver l'ancien comme attribut non hierarchique \\[2pt]
\textbf{PQ-20} & Manquants correles & Les indicateurs agricoles sont absents en bloc sur les memes communes & 21,18 \% des lignes, avec un profil de manque identique sur toutes ces colonnes & Manquant legitime : ne pas imputer. Ajouter un indicateur commune\_agricole et l'utiliser comme filtre des analyses de correlation \\[2pt]
\textbf{PQ-21} & Documentation & La definition de plusieurs indicateurs Solagro n'est pas etablie & 5 colonnes sur 46 sans definition confirmee & Confronter au guide methodologique Solagro cite dans le sujet avant de figer le dictionnaire \\[2pt]
\textbf{PQ-22} & Nommage & Convention de nommage incoherente dans le fichier source & 2 colonnes en majuscules alors que les 44 autres sont en minuscules & Renommer en ift\_bc et ift\_bc\_dif au chargement \\[2pt]
\textbf{PQ-23} & Additivite & Les mesures de variation ne sont pas additives & 52,94 \% de valeurs absentes : ces colonnes ne sont renseignees que sur le millesime 2022 & Declarer ces mesures comme semi-additives dans le schema DFM, ou les recalculer dans le DW a partir des valeurs absolues \\[2pt]
\textbf{PQ-24} & Donnee de presentation & Des codes couleur cartographiques sont melanges aux donnees & 3 colonnes, 7 a 12 modalites hexadecimales & Rejet documente. La palette de restitution est definie dans l'outil de reporting \\[2pt]
\textbf{PQ-25} & Valeurs hors bornes & Des indicateurs depassent leurs bornes logiques & p\_sau maximum 119,19 avec 25 lignes au-dessus de 100. p\_bio maximum 100,94 avec 60 lignes au-dessus de 100. ift\_t superieur a 20 sur 3 lignes (maximum 22,96) et egal a 0 sur 2 461. Population nulle sur 12 865 lignes et egale a 0 sur 18 & Regles de plausibilite a l'entree du DW : ecretage documente a 100, mise en quarantaine des outliers d'IFT, exclusion des communes a population nulle des analyses ponderees \\[2pt]
\textbf{PQ-26} & Perimetre & Des communes hors territoire national sont presentes & 99001 VENTIMILLE ITALIE sur 421 prelevements et 99004 MEIX-DEVANT-VIRTON (Belgique) sur 2 & Exclure explicitement ou creer un membre Hors France dans DIM\_GEOGRAPHIE. Tracer la decision \\[2pt]
\textbf{PQ-27} & Perimetre & Les departements d'outre-mer n'ont aucune donnee de pression agricole & 101 departements cote SISE-Eaux contre 96 departements metropolitains cote Solagro. Les codes 971 a 976 sont absents de Solagro & Exclure les DOM des besoins de correlation et le declarer, ou les conserver comme groupe temoin sans pression agricole \\[2pt]
\textbf{PQ-28} & Grain / duplication & La table de liaison commune-UDI est un instantane annuel & 51 225 groupes dupliques sur commune + reseau + quartier + debutalim apres union des 10 millesimes & Ajouter l'annee derivee du nom de fichier, puis versionner la table pont ou la dedoublonner sur le millesime le plus recent \\[2pt]
\textbf{PQ-29} & Coherence inter-fichiers & Des prelevements n'ont aucun resultat d'analyse & 2 834 678 prelevements dans DIS\_PLV contre 2 834 656 references dans DIS\_RESULT, soit 22 prelevements sans analyse & Controle de coherence automatise en fin d'ETL, avec table de rejet et seuil d'alerte \\[2pt]
\end{longtable}}

